% test-006.tex -- comando \spshort
%
% Compilar:  lualatex-dev test-006.tex
% Validar:   show-pdf-tags --xml test-006.pdf | rnv-wrapp
%            verapdf --flavour ua2 --format text test-006.pdf
%
% \spshort trabaja solo en modo texto. Las secciones A a G se generaron a
% partir de las tablas del manual de usuario: una entrada por cada
% entrada y cada alias. Cada caso es un parrafo numerado, para poder
% referirse a el al escuchar el PDF. Los comentarios "doc:" indican la
% lectura que documenta el manual; en el resto, la lectura se revisa
% escuchando. Las entradas invalidas (que detienen la compilacion) no van
% aqui.
%
% Pendiente: los alias 1.ª y 2.º de la tabla de ordinales dan error en el
% codigo, aunque el manual los lista, por eso no se incluyen. Tampoco
% 12.er, que es invalido (er solo admite terminaciones 1 y 3).
\DocumentMetadata
  {
    lang=es-CL,
    pdfversion=2.0,
    pdfstandard=ua-2,
    tagging=on,
    tagging-setup={math/setup={mathml-SE}},
  }
\documentclass{article}
\usepackage[spanish]{babel}
\usepackage{unicode-math}
\usepackage{spintent}
\usepackage{hyperref}
\hypersetup
  {
    pdfdisplaydoctitle = true,
    pdftitle           = {test-006: spshort},
  }
\begin{document}

\section*{A. Abreviaturas con letras voladas}

% doc: «doctora»
1. Entrada dr.a: \spshort{dr.a}.

% doc: «doctora»
2. Alias dr.ª: \spshort{dr.ª}.

% doc: «doctoras»
3. Entrada dr.as: \spshort{dr.as}.

% doc: «profesora»
4. Entrada prof.a: \spshort{prof.a}.

% doc: «doña»
5. Entrada d.a: \spshort{d.a}.

% doc: «doña»
6. Alias d.ª: \spshort{d.ª}.

% doc: «maría»
7. Entrada m.a: \spshort{m.a}.

% doc: «número»
8. Entrada n.o: \spshort{n.o}.

% doc: «número»
9. Alias n.º: \spshort{n.º}.

% doc: «números»
10. Entrada n.os: \spshort{n.os}.

% doc: «compañía»
11. Entrada c.ia: \spshort{c.ia}.

\section*{B. Abreviaturas comunes y expresiones latinas}

% doc: «página»
12. Entrada pág.: \spshort{pág.}.

% doc: «página»
13. Alias pag.: \spshort{pag.}.

% doc: «volumen»
14. Entrada vol.: \spshort{vol.}.

% doc: «etcétera»
15. Entrada etc.: \spshort{etc.}.

% doc: «y otros»
16. Entrada et al.: \spshort{et al.}.

% doc: «y otros»
17. Alias et. al.: \spshort{et. al.}.

% doc: «ibídem»
18. Entrada ibíd.: \spshort{ibíd.}.

% doc: «ibídem»
19. Alias ibid.: \spshort{ibid.}.

% doc: «obra citada»
20. Entrada op. cit.: \spshort{op. cit.}.

% doc: «obra citada»
21. Alias op.cit.: \spshort{op.cit.}.

% doc: «lugar citado»
22. Entrada loc. cit.: \spshort{loc. cit.}.

% doc: «lugar citado»
23. Alias loc.cit.: \spshort{loc.cit.}.

% doc: «verbigracia»
24. Entrada v. gr.: \spshort{v. gr.}.

% doc: «verbigracia»
25. Alias v.gr.: \spshort{v.gr.}.

% doc: «esto es»
26. Entrada i. e.: \spshort{i. e.}.

% doc: «esto es»
27. Alias i.e.: \spshort{i.e.}.

% doc: «por ejemplo»
28. Entrada e. g.: \spshort{e. g.}.

% doc: «por ejemplo»
29. Alias e.g.: \spshort{e.g.}.

% doc: «por ejemplo»
30. Entrada p. ej.: \spshort{p. ej.}.

% doc: «por ejemplo»
31. Alias p.ej.: \spshort{p.ej.}.

\section*{C. Títulos, profesiones y cargos}

% doc: «señor»
32. Entrada sr.: \spshort{sr.}.

% doc: «señora»
33. Entrada sra.: \spshort{sra.}.

% doc: «señorita»
34. Entrada srta.: \spshort{srta.}.

% doc: «doctor»
35. Entrada dr.: \spshort{dr.}.

% doc: «doctora»
36. Entrada dra.: \spshort{dra.}.

% doc: «doctores»
37. Entrada dres.: \spshort{dres.}.

% doc: «doctoras»
38. Entrada dras.: \spshort{dras.}.

% doc: «profesor»
39. Entrada prof.: \spshort{prof.}.

% doc: «profesora»
40. Entrada profa.: \spshort{profa.}.

% doc: «profesores»
41. Entrada profs.: \spshort{profs.}.

% doc: «ingeniero»
42. Entrada ing.: \spshort{ing.}.

% doc: «ingenieros»
43. Entrada ings.: \spshort{ings.}.

% doc: «licenciado»
44. Entrada lic.: \spshort{lic.}.

% doc: «visto bueno»
45. Entrada v. b.: \spshort{v. b.}.

% doc: «visto bueno»
46. Alias v.b.: \spshort{v.b.}.

\section*{D. Abreviaturas compuestas e institucionales}

% doc: «sociedad anónima»
47. Entrada s. a.: \spshort{s. a.}.

% doc: «sociedad anónima»
48. Alias s.a.: \spshort{s.a.}.

% doc: «estados unidos»
49. Entrada ee. uu.: \spshort{ee. uu.}.

% doc: «estados unidos»
50. Alias ee.uu.: \spshort{ee.uu.}.

% doc: «derechos humanos»
51. Entrada dd. hh.: \spshort{dd. hh.}.

% doc: «derechos humanos»
52. Alias dd.hh.: \spshort{dd.hh.}.

% doc: «juegos olímpicos»
53. Entrada jj. oo.: \spshort{jj. oo.}.

% doc: «juegos olímpicos»
54. Alias jj.oo.: \spshort{jj.oo.}.

% doc: «fuerzas armadas»
55. Entrada ff. aa.: \spshort{ff. aa.}.

% doc: «fuerzas armadas»
56. Alias ff.aa.: \spshort{ff.aa.}.

% doc: «relaciones exteriores»
57. Entrada rr. ee.: \spshort{rr. ee.}.

% doc: «relaciones exteriores»
58. Alias rr.ee.: \spshort{rr.ee.}.

% doc: «comunidades autónomas»
59. Entrada cc. aa.: \spshort{cc. aa.}.

% doc: «comunidades autónomas»
60. Alias cc.aa.: \spshort{cc.aa.}.

% doc: «posdata»
61. Entrada p. d.: \spshort{p. d.}.

% doc: «posdata»
62. Alias p.d.: \spshort{p.d.}.

% doc: «antes de Cristo»
63. Entrada a. c.: \spshort{a. c.}.

% doc: «antes de Cristo»
64. Alias a.c.: \spshort{a.c.}.

% doc: «después de Cristo»
65. Entrada d. c.: \spshort{d. c.}.

% doc: «después de Cristo»
66. Alias d.c.: \spshort{d.c.}.

% doc: «documento nacional de identidad»
67. Entrada d. n. i.: \spshort{d. n. i.}.

% doc: «documento nacional de identidad»
68. Alias d.n.i.: \spshort{d.n.i.}.

% doc: «rol único tributario»
69. Entrada r. u. t.: \spshort{r. u. t.}.

% doc: «rol único tributario»
70. Alias r.u.t.: \spshort{r.u.t.}.

% doc: «rol único nacional»
71. Entrada r. u. n.: \spshort{r. u. n.}.

% doc: «rol único nacional»
72. Alias r.u.n.: \spshort{r.u.n.}.

\section*{E. Siglas}

% doc: «documento nacional de identidad»
73. Entrada dni: \spshort{dni}.

% doc: «rol único tributario»
74. Entrada rut: \spshort{rut}.

% doc: «rol único nacional»
75. Entrada run: \spshort{run}.

% doc: «organización de las naciones unidas»
76. Entrada onu: \spshort{onu}.

% doc: «real academia española»
77. Entrada rae: \spshort{rae}.

% doc: «organización no gubernamental»
78. Entrada ong: \spshort{ong}.

% doc: «organización no gubernamental» (el código lee el plural: «organizaciones no gubernamentales»)
79. Alias ongs: \spshort{ongs}.

% doc: «unión de repúblicas socialistas soviéticas»
80. Entrada urss: \spshort{urss}.

% doc: «organización de los estados americanos»
81. Entrada oea: \spshort{oea}.

% doc: «organización mundial de la salud»
82. Entrada oms: \spshort{oms}.

% doc: «fondo monetario internacional»
83. Entrada fmi: \spshort{fmi}.

% doc: «banco interamericano de desarrollo»
84. Entrada bid: \spshort{bid}.

% doc: «organización del tratado del atlántico norte»
85. Entrada otan: \spshort{otan}.

% doc: «unión europea»
86. Entrada ue: \spshort{ue}.

% doc: «reino unido»
87. Entrada ru: \spshort{ru}.

\section*{F. Acrónimos silábicos}

% doc: «fondo de las naciones unidas para la infancia»
88. Entrada unicef: \spshort{unicef}.

% doc: «organización de las naciones unidas para la educación, la ciencia y la cultura»
89. Entrada unesco: \spshort{unesco}.

% doc: «mercado común del sur»
90. Entrada mercosur: \spshort{mercosur}.

% doc: «comisión económica para américa latina»
91. Entrada cepal: \spshort{cepal}.

% doc: «ministerio de educación»
92. Entrada mineduc: \spshort{mineduc}.

% doc: «consejo nacional de investigaciones»
93. Entrada conicet: \spshort{conicet}.

\section*{G. Números ordinales}

% doc: «primera»
94. Entrada 1.a: \spshort{1.a}.

% doc: «segundo»
95. Entrada 2.o: \spshort{2.o}.

% doc: «tercer»
96. Entrada 3.er: \spshort{3.er}.

% doc: «cuartos»
97. Entrada 4.os: \spshort{4.os}.

% doc: «quintas»
98. Entrada 5.as: \spshort{5.as}.

\section*{H. Otros ordinales y escritura en mayúsculas}

% código: «décimo»
99. Décimo: \spshort{10.o}.

% código: «undécima»
100. Undécima: \spshort{11.a}.

% código: «vigésimo primer»
101. Vigésimo primer: \spshort{21.er}.

% código: «trigésimo tercer»
102. Trigésimo tercer: \spshort{33.er}.

% código: «primeros»
103. Plural: \spshort{1.os}.

% código: «100º», sin lectura en palabras
104. Centésimo: \spshort{100.o}.

% código: «documento nacional de identidad»
105. Sigla en mayúsculas: \spshort{DNI}.

% código: «fondo de las naciones unidas para la infancia»
106. Acrónimo con mayúscula inicial: \spshort{Unicef}.

% código: «doctora»
107. Abreviatura con letra volada en mayúscula: \spshort{Dr.A}.

% código: «etcétera»
108. Abreviatura con punto final y mayúscula: \spshort{Etc.}.

\section*{I. Llave small-caps}

109. Por defecto: \spshort{dr.a}.

110. Con false: \spshort[small-caps=false]{dr.a}.

% doc: «doctora», imprime Dr. con la a en versalita
111. Letra volada con true: \spshort[small-caps=true]{dr.a}.

112. Ordinal con true: \spshort[small-caps=true]{1.o}.

113. Número con true: \spshort[small-caps=true]{n.o}.

114. Sigla de tres letras con true: \spshort[small-caps=true]{onu}.

115. Sigla de cuatro letras con true: \spshort[small-caps=true]{otan}.

116. Acrónimo de cinco o más letras con true: \spshort[small-caps=true]{unicef}.

117. Abreviatura en mayúsculas con true: \spshort[small-caps=true]{ee. uu.}.

\section*{J. Llave read-arg}

% doc: «octavo»
118. Ejemplo del manual: \spshort[read-arg={octavo}]{8.\textsuperscript{\emph{avo}}}.

119. Abreviatura nueva: \spshort[read-arg={profesor emérito}]{Prof.~em.}.

120. Sigla nueva: \spshort[read-arg={instituto nacional de estadísticas}]{INE}.

\section*{K. En contexto}

121. Varias abreviaturas en una oración: La \spshort{dra.} Pérez y el \spshort{sr.} López, de la \spshort{onu}, citaron la \spshort{pág.} 12 del \spshort{vol.} 3.

122. Ordinales en una oración: Quedó en el \spshort{3.er} lugar del \spshort{5.o} concurso.

% doc: «primero»
123. Con numerales de caja alta, como indica el manual: {\addfontfeatures{Numbers=Lining}\spshort{1.o}}.

124. Sigla y abreviatura juntas: El \spshort{rut} del \spshort{ing.} Soto figura como \spshort{n.o} 5.

\end{document}
